% THIS IS A LATEX TEMPLATE FILE FOR PAPERS INCLUDED IN THE
% *Anthology of Computers and the Humanities*. ADD THE OPTION
% 'final' WHEN CREATING THE FINAL VERSION OF THE PAPER. 
% DO NOT change the documentclass
%\documentclass[final]{anthology-ch} % for the final version
\documentclass{anthology-ch}         % for the submission

% LOAD LaTeX PACKAGES
\usepackage{booktabs}
\usepackage{graphicx}
% ADD your own packages using \usepackage{}

% TITLE OF THE SUBMISSION
% Change this to the name of your submission
\title{Not flawless, but good enough: Vision-language models and the invisible errors of machine transcription}

% AUTHOR AND AFFILIATION INFORMATION
% For each author, include a new call to the \author command, with
% the numbers in brackets indicating the associated affiliations 
% (next section) and ORCID-ID for each author.  
\author[1,2]{Author One}[
  orcid=0000-0000-0000-0000
]

\author[1]{Author Two}[
  orcid=0000-0002-8872-9474
]

% While we encourage including ORCID-IDs for all authors, you can
% include authors that do not have one by definining an empty ID.
\author[2]{Author Three}[
  orcid=
]

% There should be one call to \affiliation for each affiliation of
% the authors. Multiple affiliations can be given to each author
% and an affiliation can be given to multiple authors. 
\affiliation{1}{Some Department, Some University, Some City, Some Country}
\affiliation{2}{Another Department, Another University, Another City, Another Country}

% KEYWORDS
% Provide one or more keywords or key phrases seperated by commas
% using the following command
\keywords[eng]{computers, humanities proceedings}

% METADATA FOR THE PUBLICATION
% This will be filled in when the document is published; the values can
% be kept as their defaults when the file is submitted
\pubyear{2025}
\pubvolume{1}
\pagestart{1}
\pageend{1}
\paperorder{1}
\conferencename{Proceedings of Conference XXX}
\conferenceeditors{Editor1 Editor2}
\doi{00000/00000}  
\customhead{}  

\addbibresource{references.bib}

%%%%%%%%%%%%%%%%%%%%%%%%%%%%%%%%%%%%%%%%%%%%%%%%%%%%%%%%%%%%%%%%%%%%%%%%%%%
% HERE IS THE START OF THE TEXT
\begin{document}

\maketitle

\begin{abstract}
We argue that VLM-based transcription introduces a new form of automation bias in textual heritage work: because its errors are often fluent, idiomatic, and semantically plausible, they can feel more trustworthy than older OCR errors even when they are less faithful to the source document.

- Semantic paraphasia a language disorder where a person unintentionally substitutes a target word with another word that has a similar meaning or is in the same category.
\end{abstract}


\section{A Baskett Full of Errors}\label{sec:intro}

In 1717 John Baskett, the King's printer, published at Oxford a folio Bible of great splendor. Arthur Charlett, watching it through the press promised ``a most Magnificent English Bible,'' with a few copies on vellum for the Queen, and reported that the corrector's celebrated eye was engaged on refinements of spacing and letter-form ``invisible to vulgar Eys'' \cite[p.~195]{simpson_proof-reading_1935}. That care was spent on the appearance of the page, not on the accuracy of its text. In the headline above the chapter of Luke in which Jesus tells the parable of the vineyard, the book reads ``The parable of the vinegar,'' one misprint among many \cite[nos.~942--943]{herbert_historical_1968} (Figure~\ref{fig:fig1_vinegar}). Among historians of the English Bible, the edition holds a double reputation: superb as printing and, as text, a ``Baskett-ful of errors'' \cite[pp.~101--102]{norton_textual_2005}. Collectors came to call it the Vinegar Bible, the name under which copies have since been cataloged, sold, and exhibited.

\begin{figure}[h!]
   \centering
   \includegraphics[width=\linewidth]{fig1_vinegar_3x3.pdf}
   \caption{(1) Detail of the 1717 Vinegar Bible, whose heading to Luke~20 reads ``The parable of the vinegar.'' (2--4) Diplomatic transcriptions produced by three vision-language models; (5--6) transcriptions produced by Tesseract and Kraken. Departures from the page are marked in red. The outputs range from conspicuous glyph and segmentation errors to fluent but unsupported readings, including the replacement of \emph{vinegar} by \emph{vineyard}. William H. Scheide Library, Princeton University Library, 72.1. Photograph by the authors.}
   \label{fig:fig1_vinegar}
 \end{figure}

The Vinegar Bible offers more than a memorable, intriguing example of early modern error. It brings into view a problem that has become newly important with a new generation of automatic transcription systems: an error may now take the form of a fluent and plausible correction. We photographed the headline and asked several current vision-language models (VLMs) for a diplomatic transcription, spelling, capitalization, and punctuation exactly as printed. Figure~\ref{fig:fig1_vinegar} gives a first look at the results, transcriptions of the same photograph by three VLMs, each departing from the text in its own way (\S\ref{sec:experiments} reports the full experiment). From the clear image, most reproduced \emph{vinegar}; gpt-5.5, in the transcription shown, returned \emph{vineyard}. The substitution captures, in one word, how the errors of automatic text recognition (ATR) are changing in character. \emph{Vinegar} is the diplomatic transcription of this page; \emph{vineyard} is the emendation an editor might introduce, and mark, in an edition. The VLM output collapses that distinction: it performs a plausible editorial correction but presents the correction as documentary evidence.

The remaining outputs place this fluent correction beside more familiar transcription failures. Tesseract and Kraken preserve \emph{vinegar} but introduce local letter confusions, broken word boundaries, omissions, and garbled strings. \texttt{Llama-4-Scout}, although a VLM, behaves much like these conventional systems, repeatedly rendering long \emph{ſ} as \emph{f}. Its \emph{priefts} and \emph{faying} are visibly wrong. The two GPT outputs are smoother, but they introduce readings that are harder to
dismiss. Both write \emph{aſk} where the page prints \emph{ask}; \texttt{gpt-5-mini} gives \emph{gaue} for \emph{gave}. Whether these forms arise from misrecognition or from a model imposing what it expects eighteenth-century biblical English to look like, the documentary result is the same: the transcription records evidence that the page does not contain. Like \emph{gaue}, \emph{priefts} and \emph{faying} introduce letters absent from the source. But whereas \emph{gaue} can pass as a plausible historical spelling, \emph{priefts} and \emph{faying} visibly announce themselves as transcription errors.

This article takes that difference in error profiles as its starting point. VLMs do not replace familiar ATR failures such as omissions, glyph confusions, and broken word boundaries; they retain those possibilities while adding errors whose fluency can conceal their departure from the page. We call one such failure \emph{smooth substitution}: the unmarked replacement of an attested reading by a well-formed, contextually plausible alternative. Smooth substitutions are difficult to detect because, unlike \emph{priefts} or \emph{faying}, they may pass lexical checks, fluency-based screens, and rapid human review. They may also be directional, drawing unusual spellings, names, numbers, or phrases toward forms the model considers more probable---whether modernized, standardized, or merely more convincingly historical. Existing error rates count these differences but do not distinguish conspicuous corruption from a plausible alteration of the documentary record. We therefore develop a taxonomy of VLM transcription failures organized by their detectability and evidentiary consequences; introduce \texttt{Misnomer}, a diagnostic framework for identifying such failures; compare several VLMs with conventional ATR baselines; and test how the resulting errors behave under common quality-assurance procedures. We conclude with recommendations for evaluating and documenting VLM-derived transcriptions in cultural heritage research.

Directly beneath the misprinted headline in the parable, the chief priests ask: \emph{By what authority doeſt thou theſe things?} The question is apt for VLM transcription. When a model returns a fluent reading that the page does not support, what gives that reading its authority---and how can a user tell whether the system has transcribed the document or reconstructed it?

\section{A new kind of transcription}\label{sec:new_kind}

To know when a model has departed from the page it copies, one first has to know \emph{how} the new models read, and \emph{why} they are replacing their predecessors. For two decades, automatic transcription in the heritage sector meant dedicated recognition systems, trained per script or per collection and embedded in infrastructures such as Transkribus \cite{muehlberger_transforming_2019, nockels_understanding_2022}. The current wave is different in kind. General-purpose vision-language models are being developed for OCR, document parsing, and historical text recognition \cite{wei_general_2024, poznanski_olmocr_2025, cui_paddleocr-vl_2025, semnani_churro_2025}; benchmark studies report that, on historical documents in several languages, such models now outperform the dedicated systems themselves \cite{humphries_unlocking_2025, crosilla_benchmarking_2025}; and the surrounding discourse has begun to describe handwritten text recognition as solved, at ``expert human levels of performance'' \cite{humphries_gemini_2025}. Adoption, in short, is running ahead of evaluation.

The established systems read in stages. A conventional OCR or HTR engine segments the page into lines, recognizes character shapes or line images with a trained classifier, and decodes the result, often with the help of a lexicon, character statistics, or a language model \cite{smith_research_2018}. In that lineage, linguistic knowledge is a separable component with a bounded role: it rescores what the recognizer proposes; it does not compose the text. More recent neural recognizers narrow the distance: TrOCR pairs a pretrained image encoder with a pretrained text decoder and already writes its transcription as a sequence of wordpieces \cite{li_trocr_2022}, though the decoder is then trained to the task, and the language it carries is an aid to decoding rather than a store of expectations about what pages say. When the visual evidence fails, systems of this family degrade toward noise, the non-words and glyph confusions of Figure~\ref{fig:fig1_vinegar}~(5), the ``dirty OCR'' whose costs digital scholarship has learned to measure and to work around \cite{traub_impact_2015, alpert-abrams_machine_2016, cordell_q_2017, hill_quantifying_2019}.

The change a VLM brings is therefore not from character recognition to language-aware recognition, but in the breadth of the language brought to bear. A vision encoder feeds a large language model trained on web-scale text \cite{liu_visual_2023}, and the transcription is produced token by token, each choice conditioned on the image and on learned distributions over likely continuations. Fluency is not a by-product of this arrangement but its objective: linguistic expectation is no longer a stage applied to a reading, it is the medium in which the reading is composed. Where the signal is clear and the wording common, the two sources of evidence agree, and the results can be excellent. Where the signal weakens or the wording is rare, the prior increasingly governs the prediction. Historical collections sit disproportionately in that second condition: their spelling predates standardization, their names and numbers are unfamiliar, and their most studied features are precisely the ones a modern prior finds improbable.

Evaluation has not moved with the architecture. Transcription quality is still reported as character and word error rates \cite{neudecker_survey_2021}, edit-distance measures \cite{levenshtein_binary_1966} that price substitutions by the keystroke: \emph{rn}$\rightarrow$\emph{m}, an obvious glyph confusion, costs two edits, and \emph{vinegar}$\rightarrow$\emph{vineyard} costs the same two. For classical OCR this pricing was reasonable: its errors were noise, visible and roughly interchangeable, so counting them was a fair measure of the damage. The same arithmetic misleads for a system whose characteristic error is a real word in a plausible sentence. A rising error rate still signals trouble, but a falling one no longer certifies fidelity, because the errors that survive are precisely those that read like correct transcription and can only be caught by comparing the output against the image.

The field of automatic speech recognition (ASR) faced this evaluation problem first. In 2003, its researchers were asking whether word error rate predicts understanding at all \cite{wang_is_2003}. A transcript that drops a filler word, an ``uh'' or a ``like,'' takes a full penalty on its error rate even though the meaning survives. A transcript that drops a ``not,'' or hears ``sell'' as ``buy,'' pays about the same penalty and reverses what the speaker meant. Newer systems sharpened the problem. Older recognizers matched sounds to words and used a language model only to rank candidate readings; today's speech models generate the transcript directly, as a language model generates any text, and their errors come out fluent: confident, well-formed transcripts of things nobody said \cite{koudounas_hallucination_2025}. The field's answer was to measure meaning, scoring a hypothesis by how far it drifts from the sense of the utterance \cite{kim_semantic_2021}. Documentary transcription needs a different measure. A metric of meaning has no use for the difference between \emph{gaue} and \emph{gave}, and it can say of \emph{vineyard} for \emph{vinegar} only that two related words were exchanged, nothing about which of them the page prints or what hangs on that.

Our taxonomy therefore asks two questions about each failure. Can the error be recognized from the transcription itself? And if the reading is accepted, what documentary evidence has been lost, altered, or invented? The next section organizes the observed failure modes by these two properties: their detectability and their evidentiary consequence.

\section{A taxonomy of fluent failure}\label{sec:taxonomy}

The five machine transcriptions in Figure~\ref{fig:fig1_vinegar} contain, between them, most of the error types that automatic transcription can produce. This section sorts them into a taxonomy of six kinds of error, several of which are new, becasue they depend on the language capabilities that VLMs bring to the task. About each kind we ask two questions. Can the error be recognized from the transcription alone, by a spelling check, a rule, a stated convention, or does recognizing it require comparison with the image? And if the error goes unrecognized, what happens to the documentary evidence: is something lost, altered, or invented? The first question decides how much verification a reading needs; the second, what is at stake when that verification does not happen. We define the six
kinds below. The experiments in \S\ref{sec:experiments} then measure how often each occurs across systems and material.

\begin{enumerate}
    \item The most familiar kind is \textsc{\textbf{Garble}}: output that is not a word, or that breaks and joins words across their boundaries. \emph{Priefts} and \emph{faying} in Figure~\ref{fig:fig1_vinegar} belong here, as do most of errors that Tesseract generated. This is the error class digital scholarship knows as ``dirty OCR'', whose costs for search and analysis have been measured at length \cite{traub_impact_2015, hill_quantifying_2019}. Every system in our comparison below produces garble, and conventional engines produce little else when the image defeats them. Its redeeming property is that it is self-announcing: a spelling check flags it, a reader stumbles over it, and the stumble prompts exactly the verification the passage needs.
    
    \item \textsc{\textbf{Silent Normalization}} changes the spelling of words without changing their semantics. In Figure~\ref{fig:fig1_vinegar}, Kraken renders each long~\emph{ſ} as \emph{s}. Editors know the layer being rewritten as the `accidentals' of a text, usually arbitrary, frequently altered by printers, typesetters, or scribes during production rather than by the author \cite{greg_rationale_1950}. An edition that regularizes spelling is said to normalize silently, with the difference that the edition states its policy in the front matter. How much policy a transcription system states varies sharply between the two families. The Kraken model in Figure~\ref{fig:fig1_vinegar} is distributed via Zenodo with a description of its transcription convention \cite{gabay_catmus-print_2024}, and its model file lists its complete output alphabet: 220 characters, among which \emph{ſ} does not appear. Every long~\emph{ſ} it reads must therefore surface as some other character. Both facts, and the discrepancy between them, can be verified, because the inventory is inspectable and the convention is published. A VLM publishes neither. Its tokenizer operates on bytes, so any Unicode string is a possible transcription \cite{kudo_sentencepiece_2018, radford_language_2019}, and the convention it applies is whichever spelling its training data made most probable (stated nowhere). A normalization can only be detected against a declared convention, and only one of the two families provides a declaration.

    \item \textsc{\textbf{Fake Authenticity}} runs in the opposite direction: the model makes the text look older than the page. The GPT models write \emph{aſk} where the type reads \emph{ask}, and \emph{gaue} where it reads \emph{gave}. Linguists call this pattern hypercorrection, the over-application of a register a writer is imitating; here the register is age. Producing it requires an idea of what period text should look like. Conventional recognizers have no such idea, and for the Kraken model discussed above the direction is closed outright, since \emph{ſ} is absent from its alphabet; in our runs, fake authenticity appears only in VLM output. Taken together with silent normalization, it also corrects a tempting simplification. The prior does not pull the text toward the present; it pulls toward whatever the model expects, and on an eighteenth-century page it expects old spelling. The cost falls on studies that date or localize documents by their orthography, and histories of spelling change. Moreover, a fabricated archaism, once transcribed and published, can also enter the training data of later models, where it counts as an attested period form.

    \item \textsc{\textbf{Entity Smoothing}} replaces proper names and rare terms with more familiar ones. Names invite this failure more than any other part of a text. They form an open class, so no dictionary can vouch for them, and they are among the least frequent words on any page, so the model's expectations give them the least support. The two system families fail differently here. A conventional engine mangles an unfamiliar name into garble, which is visible. A VLM replaces it, sometimes with a different well-formed name, sometimes with a common word. Medical speech recognition has documented the same class and built entity-level metrics for it, because a transcript can score well overall (e.g., a low CER) while getting exactly the clinically decisive terms wrong \cite{afonja_performant_2024}. In historical collections the stakes are retrieval and identity: a smoothed name cannot be searched for, linked to an authority record, or counted. The names most exposed to this problem are often also the rarest, and rare names belong disproportionately to the people and places that archives document least.

    \item \textsc{\textbf{Confabulated Continuation}} operates on longer stretches. Where the visual signal fails entirely, a generative model can produce a fluent passage with no source on the page, the transcription counterpart of hallucination in speech recognition and summarization \cite{koudounas_hallucination_2025, maynez_faithfulness_2020}. A conventional engine in the same situation produces garble or nothing. Because the invented passage is grammatical throughout, no check that operates on single words will object to it.

    \item The last kind is the \textsc{\textbf{Smooth Substitution}}: a real word, plausible in its sentence, in place of what the page attests. \emph{Vineyard} for \emph{vinegar}, the substitution we began with in \S\ref{sec:intro}, shows every property of the class. The replacement is a correctly spelled word, and it fits its context: the misprinted heading stands directly above the chapter in which Jesus tells the parable of the vineyard, so \emph{vineyard} is exactly the word that the surrounding page leads a reader, or a model, to expect there. Nothing in the transcription signals that anything went wrong; the error can only be established by holding the output against the image. Substitutions of this kind also tend to have a direction: the replacement is usually the more expected word where the page has something rarer, whether a modern form, a standard spelling, or a familiar name.\footnote{Numbers are the purest form of the problem. Any sequence of digits is well formed, so a wrong year, 1717 in place of 1771, passes every check that operates on the text itself; only an external source or the image can settle it.} The failure resembles semantic paraphasia, the speech disorder in which a person produces a word related to the one intended, \emph{cat} for \emph{dog}, without noticing the exchange \cite{goodglass_understanding_1993}. Adjacent fields know this class and agree on its difficulty. Spelling correction distinguishes non-word errors, which a dictionary catches, from real-word errors, where the mistake is itself a dictionary word; detecting real-word errors was ranked the hardest problem of that field, and the standard methods all find them the same way, through the surrounding words \cite{kukich_techniques_1992, mays_context_1991, golding_winnow-based_1999}. A smooth substitution defeats that method too. Because the replacement is produced by a model of context, it agrees with the surrounding words as well, and the one place where a real-word error could still be seen shows nothing.\footnote{Human readers fare no better. Asked ``How many animals of each kind did Moses take on the Ark?'', about half of experimental participants answer ``two,'' without noticing that the Ark was Noah's; readers who demonstrably know the story fail at comparable rates, and the illusion strengthens with the semantic relatedness of the substituted name \cite{erickson_words_1981, van_oostendorp_moses_1990}.} Textual criticism documented the scribal counterpart under the name \textit{banalization}: copyists tended to replace rare readings with common ones, a drift so regular that editors made a rule of it (\textit{lectio difficilior potior}): where manuscript witnesses disagree, prefer the more difficult reading \cite{timpanaro_freudian_1976}. Conventional ATR systems and VLMs can both produce a plausible wrong word; how often each does, and in which direction its replacements lean, is what our experiments measure.

\end{enumerate}

\section{Misnomer: measuring what the metrics miss}\label{sec:misnomer}

Character and word error rates remain necessary infrastructure: they are computationally cheap, comparable, and universally understood. But they are inherited from an earlier generation of systems, and they miss precisely the errors that VLM-aided transcription adds, the smooth substitutions. This section describes the instrument we built to catch what they miss. We call it \texttt{Misnomer}, an open-source evaluation framework that compares a transcription with a verified reference, finds every place where the two disagree, and classifies each disagreement. 

To detect a transcription error, one compares the machine's output with the verified reference: wherever the two differ, an error has been found. The difficulty lies in deciding what kind of difference each one is. The pipeline we propose is a cascade of checks, ordered from the most mechanical to the least. Empty and refused responses are set aside first and counted as failures of coverage, as are lines whose content as a whole no longer corresponds to the reference. The remaining text is aligned with the reference word by word; an aligned pair whose words differ is a substitution,\footnote{Differences of case and of surrounding punctuation are disregarded at this stage, so a substitution is always a difference in the word itself.} and substitutions are where the kinds of \S\ref{sec:taxonomy} occur. Each substitution then meets four checks. (1)~A dictionary identifies garble: words that appear in no lexicon. (2)~Numbers, and any string containing a digit, form a class of their own. If a transcription reads 1775 where the page reads 1777, nothing in the text reveals the error: 1775 is not misspelled, no dictionary lists years, and the surrounding sentence accepts either date equally well. The pipeline therefore sets digit-bearing substitutions apart, so that they are counted, and can be flagged for checking against an external source, rather than passed along as if they were words. (3)~A transcription convention, a machine-readable statement of the editorial practices in force, selected or defined by the user, identifies pairs that are the same text rendered under different standards (for instance the \emph{u}/\emph{v} interchange of early modern print): modernized spelling in one direction, fabricated period spelling in the other. (4)~Character-level alignment separates misread words from words broken or joined at their boundaries. What survives every check is the residue: a real word, correctly spelled, fitting its context, and different from what the page attests. For that residue the framework computes one further quantity, defined below: which of the two words a language model would sooner expect at that position, the word the system transcribed or the word the reference attests, and how strong that preference is.

To make this comparison for a substitution that has passed all four checks, we need three things. We need the word the model wrote, $\hat{w}_i$. We need the word the page actually contains, $w^{*}_i$. And we need the context in which the two words compete, $\mathbf{c}_i$: the words that come before position $i$. All three are already in hand from the word-by-word alignment of transcription and reference with which the pipeline began.\footnote{For the context we use the words of the reference, not the words of the transcription, so that errors elsewhere in the line cannot influence the test at this position.} The question we put to each triple concerns the direction of the error: of the two words, which one does the context make the easier, more expected reading? Nothing examined so far can answer this, because at this stage both words are correctly spelled and both fit the sentence.

Asking that question of a machine requires a machine that has expectations. A causal language model is exactly that: given a sequence of words, it assigns a probability to every possible continuation. We therefore employ one as an \emph{auditor}; in our experiments, Qwen2.5 with half a billion parameters, pinned to a fixed version.\footnote{The choice follows three considerations. The auditor can be small, because its role calls for ordinary linguistic expectation rather than knowledge or reasoning. It should be open and version-pinned, so that its judgments can be reproduced exactly. And it must be independent of every system under evaluation, so that no model grades its own output. Any causal language model can serve in the role, and the choice is a declared parameter of the framework; a different auditor brings a different linguistic prior, which is why \S\ref{sec:experiments} repeats the analysis under a second, unrelated auditor.} For each substitution, the auditor reads the context and scores both candidates:
%
\begin{equation}
    \Delta_i \;=\; \log P_{\theta}\!\left(\hat{w}_i \mid \mathbf{c}_i\right)
             \;-\; \log P_{\theta}\!\left(w^{*}_i \mid \mathbf{c}_i\right).
    \label{eq:delta}
\end{equation}
%

The two terms score the transcription's word and the reference word in the same context; the difference is the auditor's preference between them. A positive $\Delta_i$ means the auditor finds the transcription's word the more expected one: the substitution moved the text toward what the context predicts. A negative $\Delta_i$ means the reference word was the more expected one: the substitution moved the text away from what the context predicts. Substitutions that originate in the image usually score this way, because a system that confuses similar letterforms (\emph{rn} for \emph{m}, \emph{e} for \emph{o}) produces the word the shapes suggest, and that word is rarely the one the context predicts.

To illustrate the procedure, suppose the reference reads ``the ancient
festival was reviled in the years that followed,'' and the machine transcription renders \emph{reviled} as \emph{revived}. Both words are real and correctly spelled, and the transcribed sentence, though it contains an error, is grammatical. Conditioned on the preceding words, ``the ancient festival was,'' the auditor assigns \emph{revived} the higher probability: $\Delta = +4.83$.\footnote{Had the reference read \emph{revived} and the transcription \emph{reviled}, the score would be $-4.83$. The procedure treats the two candidates identically, so exchanging them changes only the sign, and equal words score zero.} The transcription has replaced the reference word with the word a language model expects at that position. This is the smooth substitution of Section~\ref{sec:taxonomy}: the error yields a real word in a sentence that reads correctly, so the transcription itself carries no signal of the exchange, and it becomes visible only in comparison with the reference.

Run over a corpus, the pipeline assigns every substitution a category, and three numbers summarize what it finds for each system. The invisible error rate (IER) counts the substitutions that survive every filter --- real words, correctly spelled, consistent with the declared convention, whose characters correspond to one word in the reference --- per thousand words of reference text. Everything in this count is deterministic. The auditor enters only afterward and contributes the second number, the smooth-substitution share: the fraction of a system's invisible errors with $\Delta_i > 0$, the substitutions that exchanged the reference word for the word the context made more expected. The third number is counted over lines rather than single words. A line is clean-looking when its transcription shows no visible fault: no garbled word, no misspelling, nothing against the declared convention. The false-clean rate is the fraction of clean-looking lines that nevertheless contain at least one invisible substitution --- the lines a reader without the original would accept. Section~\ref{sec:experiments} reports all three for every system we test, alongside the conventional error rates.

\section{Systems and material}
\label{sec:design}

To measure how often working transcription systems produce smooth substitutions, we run ten of them over the same corpus and score every output with the pipeline of Section~\ref{sec:misnomer}. The ten differ chiefly in how much knowledge of language they bring to the image. Two are conventional engines with no learned language model: Tesseract \cite{smith_overview_2007} and Kraken \cite{kiessling_kraken_2019} with the McCATMuS model \cite{chague_mccatmus_2024}. Two are runs of the same handwriting recognizer, PyLaia \cite{puigcerver_are_2017} trained on IAM, decoded once without and once with its six-gram language model \cite{tarride_improving_2024}: a pair that differs in nothing except the strength of its linguistic prior. Six are vision--language models: GPT-4o, GPT-5, Gemini~3~Pro, and Claude Opus~4.6, accessed through their APIs,\footnote{Snapshot identifiers and access dates for the API-served models accompany the released data.} and two with open weights, Qwen3-VL \cite{qwen_team_qwen3-vl-8b-instruct_2025} and Chandra \cite{datalab_chandra_2025}, the latter specialized for document transcription.

The corpus we use is the IAM handwriting database \cite{marti_iam-database_2002}: 10,373 line images in which several hundred writers copied printed modern English prose. The hands are contemporary, the orthography is modern, and the scans are clean; the difficulties that explain transcription errors on historical material --- unfamiliar scripts, distant spelling conventions, degraded pages --- are absent, so an error here cannot be attributed to the material. The dataset is also familiar: IAM has been the standard handwriting benchmark for two decades, so we can assume every model we test has seen it in training. Benchmark studies treat such contamination as a flaw. For our question it is useful: familiarity should suppress substitution, so whatever a system substitutes here, it can be expected to substitute at least as often on material it has never seen.

Every system receives the same 10,373 images, and the vision--language models receive the same short transcription instruction. All outputs pass through the pipeline of Section~\ref{sec:misnomer}, with one dictionary and one auditor throughout. Responses that come back empty, refused, or unrelated to the image count as coverage failures, reported alongside the error rates rather than inside them. PyLaia is scored on IAM's test split only, because it was trained on the training split. Confidence intervals are bootstrapped over lines.

[[[4,500 words up until this part! 1,500 for the results and conclusion!]]]

\section{Experiment}\label{sec:experiments}

The ten systems were asked for 103,730 line transcriptions in all. The first result concerns the requests that produced no transcription to score. GPT-5 returned an empty response for 2,285 of its 10,373 lines, 22 percent of the corpus; Tesseract did so for 913 lines, 9 percent. Gemini answered every line, but 74 of its responses contained the model's reasoning about the image in place of a transcription of it; a length filter catches these, and they too count as coverage failures. GPT-4o left 2 percent of lines unanswered, and the remaining six systems returned output for every line. Published comparisons rarely report these numbers, yet every error rate silently conditions on them: a system is scored only on the lines it answered, so GPT-5's rates in what follows describe the 78 percent of the corpus it transcribed. Table~\ref{tab:census} therefore reports coverage as a column of its own, beside the error rates, for all ten systems.

\begin{table}[t]
\centering
\begin{tabular}{l r r r r r}
\toprule
System & CER\,\% & WER\,\% & IER & Smooth\,\% & False-clean\,\% ($n$) \\
\midrule
Tesseract & 37.0 & 70.4 & 170.0 & 7 & 55 (58) \\
Kraken (McCATMuS) & 25.0 & 64.4 & 124.0 & 7 & 66 (58) \\
\midrule
PyLaia\textsuperscript{t} & 7.0 & 22.7 & 35.6 & 12 & 14 (724) \\
PyLaia + LM\textsuperscript{t} & 7.2 & 21.3 & 43.1 & 15 & 20 (890) \\
\midrule
GPT-4o & 6.1 & 6.9 & 28.5 & 46 & 17 (8,979) \\
Chandra & 5.3 & 5.7 & 18.1 & 53 & 12 (9,064) \\
Qwen3-VL & 4.8 & 5.2 & 16.7 & 45 & 11 (9,251) \\
GPT-5 & 4.1 & 3.9 & 10.1 & 53 & 8 (7,414) \\
Claude Opus 4.6 & 4.3 & 3.8 & 8.9 & 41 & 7 (9,602) \\
Gemini 3 Pro & 4.3 & 3.6 & 6.7 & 45 & 5 (9,530) \\
\bottomrule
\end{tabular}
\caption{The three measures of Section~\ref{sec:misnomer}, beside the conventional error rates, for all ten systems over the 10,373 IAM lines. CER and WER are micro-averages over the lines a system answered; GPT-5 answered 78\% of lines, Tesseract 91\%, GPT-4o 98\%, all others 100\%. IER is invisible substitutions per thousand reference words. Smooth\,\% is the smooth-substitution share after the precision screens. False-clean\,\% is the share of clean-looking lines containing at least one invisible substitution, with the number of clean-looking lines in parentheses. The two PyLaia rows, marked \textsuperscript{t}, are scored on IAM's test split only, because that model was trained on the training split. Digit-bearing substitutions are diverted by the numeric gate and cannot appear as invisible substitutions; bootstrap confidence intervals and unscreened values accompany the released data.}\label{tab:census}
\end{table}

\subsection{Towards Semantic Error Detection in Vision--Language Models}

== this becomes part of smooth substitution 6. ===
We define a semantic error as a class of machine errors in Automatic Text Recognition (ATR) in which the model transcribes text that is linguistically plausible but unfaithful to the source image because visual evidence is overridden by the language priors of a pre-trained language model. Unlike classical recognition errors, which stem from degraded, ambiguous, or out-of-distribution visual signals, semantic errors occur even when the decoder emits what is probable under its learned language distribution rather than what is actually written. Because these priors are intrinsic to the generative language-modeling objective on which such systems are built, semantic errors are not incidental artifacts but a systematic consequence of using a generative language model for transcription. 

===

Our method measures semantic error in the output of any vision--language model (VLM) by comparing a predicted transcription against a ground-truth transcription of the same document image. The central design choice is a \emph{model-agnostic} scoring pipeline. Given that some models expose their internal language model and others do not, we decided to score only the outputs. 

All scoring is performed by a fixed, versioned set of internal models that are independent of the transcription system. This makes scores reproducible and comparable across systems, collections, and institutions, and it lets a single scorer audit output from a local VLM, a remote API, or a classical OCR engine on equal terms.

\subsection{Word alignment}

Predicted and ground-truth strings are first tokenized into words. When a
transformer tokenizer is available, subword tokens are grouped into words using
their character offsets---robust across byte-level BPE (Qwen2.5, GPT-2),
SentencePiece, and WordPiece tokenizers---after which leading and trailing
punctuation is stripped and punctuation-only tokens are discarded. The two word
sequences are aligned by minimum edit distance using the same dynamic-program that underlies word error rate (WER); matching is case-insensitive. The
alignment labels each position as a \textsc{match}, \textsc{substitution},
\textsc{insertion}, or \textsc{deletion}. Only substitutions can carry a
semantic error, so only substituted word pairs receive a composite error score;
insertions, deletions, and matches are scored zero and instead feed the
aggregate rate statistics described below.

\subsection{Two signals}


Each substituted pair---a predicted word $\hat{w}$ aligned to a ground-truth
word $w^{*}$---is characterized by two independent signals.

\paragraph{Signal 1: surprisal contrast between the correct and the predicted word.}
Using the internal language model, we compare how surprising the ground-truth
word $w^{*}$ is against how surprising the word the transcription model actually
produced, $\hat{w}$, both evaluated in the \emph{same} ground-truth left-context.
Writing $s(w)=-\log P(w \mid \text{context})$ for the contextual surprisal of a
word, the signal is the signed contrast
\[
  \Delta \;=\; s(\hat{w}) - s(w^{*}).
\]
A negative contrast ($\Delta<0$) means the model finds the \emph{misread} word
more fluent than the truth: the page says ``vinegar,'' but the model has learned
``vineyard'' during training and prefers it here. This is exactly the condition
under which a semantic substitution is most tempting and most damaging.

At first, we used the absolute surprisal of the correct
word alone. However, that conflated two different situations: a genuine override and a
rare-but-correct word such as a name, a place, an unusual term, that is simply
surprising wherever it appears. Because the correct word is often the rarer one,
absolute surprisal systematically assigned high scores to correct entities. The
contrast removes this confounding feature: a shared rarity term cancels, so a transcribed rare word correctly yields $\Delta \approx 0$ regardless of how
surprising it is, and the signal rises only when the prior language actively favors the wrong word.

To obtain a stable per-word value we evaluate the surprisal of each candidate on
only the \emph{first} subword token of its span. Continuation tokens within a
multi-token word are conditioned on the word having already started; they are
near-certain predictions whose log-probabilities are close to zero and would, if
averaged in, systematically deflate the surprisal of longer words. The first
token is the one the model predicts purely from context, so it carries the
genuine surprisal signal. We map the contrast to an override risk in $[0,1]$ with
a temperature-scaled logistic,
\[
  r \;=\; \sigma\!\left(-\Delta / T\right),
\]
which at $T=1$ is exactly the model's relative preference for the misread word,
$P(\hat{w}\mid\text{context}) / \bigl(P(\hat{w}\mid\text{context}) +
P(w^{*}\mid\text{context})\bigr)$. This risk replaces the per-word perplexity
term in the composite score; the earlier absolute-surprisal behavior remains
available as a configuration option.

\paragraph{Signal 2: semantic proximity of the substitution.}
Using a fixed internal sentence-transformer, we measure how close in meaning the
predicted word is to the correct word. Each word is embedded and the two
embeddings are compared by cosine similarity, so a near-synonym yields a high
similarity and an unrelated word a low one. Because bare single-word inputs
produce unreliable embeddings each word is first
wrapped in a minimal sentence frame, ``\texttt{A \{word\}.}'', before encoding.

The counterintuitive but essential point is that \emph{high} similarity marks
the dangerous case. A substitution that means almost the same thing as the
original reads naturally and survives human review; a substitution that means
something unrelated is caught on sight. We therefore use this signal both to
grade severity and to classify the substitution: pairs with similarity at or
above a threshold (default $0.35$ for frame-encoded word similarity with the
default embedder) are labeled \textsc{semantic} errors---the plausible,
review-surviving class this work targets---and pairs below the threshold are
labeled \textsc{obvious} errors, the self-evident mistakes that classical
metrics already surface.


\subsection{Composite word score}

The two signals are combined into a single severity score per substituted pair:
%
\begin{equation}
    s_i \;=\; w_{\mathrm{ppl}}\cdot \tilde{p}(w^{*}_i \mid \mathbf{c}_i)
            \;+\; w_{\mathrm{sem}}\cdot \mathrm{sim}(\hat{w}_i,\, w^{*}_i),
    \label{eq:composite}
\end{equation}
%
where $\tilde{p}\in[0,1]$ is the log-scale-normalized perplexity of the correct
word in context (Signal 1), $\mathrm{sim}\in[0,1]$ is the cosine similarity
between the predicted and correct word embeddings (Signal 2), and the score is
clipped to $[0,1]$. The default weights are $w_{\mathrm{ppl}}=0.4$ and
$w_{\mathrm{sem}}=0.6$, placing slightly more emphasis on semantic proximity, on
the argument that a plausible-meaning substitution is the more consequential
failure. A high $s_i$ therefore identifies the most dangerous substitutions:
those that are semantically close to the truth \emph{and} occur where the
correct word was contextually surprising---replacements that read more smoothly
than the source document itself and are the most likely to pass unnoticed
through careful human review.


\begin{table}[h]
\centering
\caption{Four real substitutions from the corpus, scored by embedding
similarity and override risk (the probability the reference LM prefers the
predicted word over gold in the same context). Similarity alone decides
SEMANTIC vs.\ OBVIOUS; override risk is an independent signal about
fluency, not a second vote on the category -- rows 1 and 3 share nearly
identical override risk (0.99 vs.\ 0.97) yet land in opposite categories
because similarity differs.}
\label{tab:signal-examples}
\begin{tabular}{rlrrl}
\hline
Substitution (gt $\rightarrow$ pred) & Similarity & Override risk & Category \\
\hline
disarmament $\rightarrow$ answer & 0.36 & 0.99 & SEMANTIC \\
good $\rightarrow$ fine & 0.60 & 0.00 & SEMANTIC \\
rallied $\rightarrow$ raised & 0.29 & 0.97 & OBVIOUS \\
Gaitskell $\rightarrow$ Cathell & 0.42 & 0.05 & GARBLE \\
\hline
\end{tabular}
\end{table}

\subsection{Experiment 1 - Vinegar Bible}

\subsection{Experiment 2 - IAM}
Experiment 2 evaluates semantic error on the \href{https://huggingface.co/datasets/Teklia/IAM-line}{IAM handwriting corpus}. Each row in our evaluation outputs corresponds to one IAM line image paired with its ground-truth transcription. Given that the handwritten examples are in English and that the IAM dataset is very likely part of a model's training data, the presence of semantic errors suggests that this type of error is inherent to the transformer architecture. We used the same prompt and images for all models.  

\begin{table}[h]
\centering
\caption{Semantic errors are common across all evaluated VLM OCR systems.}
\label{tab:semantic-errors-common}
\begin{tabular}{lrrrr}
\hline
Model & IAM Lines & Lines with $\geq 1$ semantic error & Percent (\%) & Max errors/line \\
\hline
ChatGPT-4o             & 10,373 & 1,932 & 18.6 & 7 \\
ChatGPT-5              & 10,373 & 1,645 & 15.9 & 5 \\
Chandra                & 10,373 & 1,447 & 13.9 & 5 \\
Qwen3-VL-8B-Instruct   & 10,373 & 1,391 & 13.4 & 5 \\
Claude Opus 4.6          & 10,373 & 965 & 9.3 & 3 \\
Gemini 3.0       & 10,373 & 825 & 8.0 & 10 \\
\hline
All VLM systems (pooled) & 62,238 & 8,205 & 13.2 & 10 \\
\hline
\end{tabular}
\end{table}
The pattern is consistent across all evaluated VLM systems: semantic errors are not edge cases but a regular machine-error mode. Depending on model, 8.0\% to 18.6\% of IAM lines contain at least one semantic error, with a pooled rate of 13.2\% (8,205 of 62,238 VLM line outputs). In other words, roughly one in eight VLM transcriptions in this experiment carries meaning-level risk despite remaining fluent and readable.Taken together, the IAM results support our central claim: semantic error is a distinct and measurable VLM transcription failure type that is only partially captured by conventional surface-form metrics such as CER and WER.



\subsection{Experiment 3 - Donkeys}

\subsection{Experiment 4 - Churro}

So far, we have investigated semantic error in English-language documents.
For computational humanities scholars, materials are often in multiple languages and scripts.
To get some bearing on how semantic error varies across languages, scripts and historical text, we identified the  Churro dataset from Stanford OVAL as a suitable testbed.

The Churro 3B model is a continued-pretraining of
Qwen2.5-VL-3B-Instruct trained specifically to transcribe historical
documents; the \texttt{churro-dataset} test split on which we evaluate spans
29 language clusters and a comparable range of scripts, from Latin and
Cyrillic to Han, Devanagari, and Khmer. Historical texts are particularly relevant to this study, because they contain archaic word forms and characters that are infrequent in modern web-scale training data. The Churro experiments focus on this problem of modern priors and demonstrates that continued pretraining on a corpus of historical documents provides the most accurate models for research. The Churro data and model allow us to evaluate the effects of continued pretraining on semantic error. Is continued pretraining an effective way to reduce semantic error rates?

We compare Churro 3B against its own base
model, Qwen2.5-VL-3B-Instruct, on identical pages, so that every difference
we report below is attributable to Stanford's continued pretraining recipe specifically,
and not to a change in model architecture or scale.

We compute our semantic error score alongside CER and WER, the metrics
reported for Churro in the original paper, over
the same aligned ground-truth region so that both families of metrics
describe the same transcribed text. The two disagree more often than they
agree. Standard edit-distance metrics are, first, poorly behaved on several
of the scripts in the test set: because CER and WER are unbounded when a
prediction runs longer than its reference, character error rates exceed
1,000\% for Chinese, Japanese, Sanskrit, and Khmer, figures that describe a
numerical blow-up rather than a level of transcription quality. Second, and
more consequentially, low CER does not imply semantic fidelity: 318 pages in
our combined evaluation (13.6\%) have a CER under 10\% yet more than five
semantic errors each, split almost evenly between the base and fine-tuned
models. These are transcriptions that a CER-only evaluation would certify as
successful while nevertheless containing the class of fluent, meaning-altering
substitution this paper is concerned with, and fine-tuning does not change the
rate at which this occurs---whatever Stanford's recipe optimizes for, it is
not the gap between surface accuracy and semantic faithfulness.

Considered on its own terms, Churro fine-tuning does reduce semantic
error substantially: averaged over all 29 languages, the semantic error rate
falls from 0.270 to 0.130 errors per ground-truth word
(Table~\ref{tab:churro-delta}), roughly a halving. But the average obscures
more than it reveals. The improvement ranges from a near-halving of semantic
error in Turkish and Persian to negligible effects in languages such as
English, and it is entangled with a second, less visible shift: the
fine-tuned model transcribes less of the page. Page coverage, the fraction
of the ground-truth text a model actually attempts, falls for 28 of the 29
languages under fine-tuning (Figure~\ref{fig:churro-coverage}), by up to 
49 points for Italian. Part of Churro's apparent gain in semantic fidelity is
therefore a gain in caution: transcribing a shorter, more confident span of
the page rather than the same span more faithfully. A fine-tuning result
reported as a single aggregate error rate can, in other words, be a coverage
effect in disguise.

Khmer is the one language in the test set where this pattern inverts, and it
is worth a brief note because the reversal is qualitative rather than
marginal. The base model does not produce Khmer script for these pages at
all; it hallucinates fluently, repeating Chinese text unrelated to the source,
consistent with Khmer's negligible presence in its pretraining data.
Churro's fine-tuning does teach the model to emit genuine Khmer script, but
on the same pages the fine-tuned model falls into a repetition loop,
generating a short Khmer phrase far past the length of the source page. The
result inflates edit-distance error, and our semantic error count together,
and it inflates measured coverage as well, since a length-based coverage
heuristic cannot distinguish a longer correct transcription from a longer
degenerate one. Khmer therefore looks, in the aggregate tables, like
continued pretraining's single failure case; read against the actual output, it is
better described as a real capability gain---the model now attempts the
correct script---compromised by a decoding failure that our metrics,
semantic and traditional alike, register only as a much worse number.

\begin{table}[h]
  \centering
  \caption{Semantic errors per ground-truth word, base model vs.\ Churro 3B,
    for all 29 languages in the test split, ordered by mean semantic error
    rate across both models. $\Delta = \text{Churro} - \text{Qwen}$;
    negative values indicate fine-tuning reduced semantic error.}
  \label{tab:churro-delta}
  \begin{tabular}{lrrr}
    \toprule
    Language & Qwen 2.5-VL-3B & Churro 3B & $\Delta$ \\
    \midrule
    Norwegian  & 0.138 & 0.022 & $-0.116$ \\
    Portuguese & 0.123 & 0.031 & $-0.092$ \\
    Spanish    & 0.113 & 0.031 & $-0.082$ \\
    Italian    & 0.064 & 0.021 & $-0.044$ \\
    German     & 0.058 & 0.014 & $-0.043$ \\
    Latin      & 0.051 & 0.019 & $-0.032$ \\
    French     & 0.055 & 0.009 & $-0.046$ \\
    Dutch      & 0.051 & 0.007 & $-0.044$ \\
    Catalan    & 0.045 & 0.013 & $-0.033$ \\
    Swedish    & 0.047 & 0.009 & $-0.038$ \\
    English    & 0.043 & 0.013 & $-0.030$ \\
    Romanian   & 0.035 & 0.002 & $-0.034$ \\
    Polish     & 0.011 & 0.015 & $+0.004$ \\
    Finnish    & 0.012 & 0.006 & $-0.006$ \\
    Slovenian  & 0.008 & 0.002 & $-0.006$ \\
    Czech      & 0.005 & 0.003 & $-0.002$ \\
    Arabic     & 0.000 & 0.003 & $+0.003$ \\
    Persian    & 0.000 & 0.003 & $+0.003$ \\
    Chinese    & 0.002 & 0.000 & $-0.002$ \\
    Bangla     & 0.001 & 0.001 & $+0.000$ \\
    Greek      & 0.000 & 0.001 & $+0.001$ \\
    Japanese   & 0.001 & 0.000 & $-0.001$ \\
    Khmer      & 0.001 & 0.000 & $-0.001$ \\
    Hebrew     & 0.000 & 0.000 & $+0.000$ \\
    Bulgarian  & 0.000 & 0.000 & $+0.000$ \\
    Hindi      & 0.000 & 0.000 & $+0.000$ \\
    Sanskrit   & 0.000 & 0.000 & $+0.000$ \\
    Turkish    & 0.000 & 0.000 & $+0.000$ \\
    Vietnamese & 0.000 & 0.000 & $+0.000$ \\
    \midrule
    \textbf{Mean (29 languages)} & \textbf{0.030} & \textbf{0.008} & \textbf{$-0.022$} \\
    \bottomrule
  \end{tabular}
\end{table}


\section{``Good enough for what?'' Guidelines for LIS and humanities}

errors will continue to exist - neede is a way to recognize it, quantify it, ten if needed correct it - but what -- 


\section{Future directions}
The present framework requires a verified ground-truth transcription against which the predicted text is aligned. Alignment is performed at the word level using edit-distance-based sequence alignment analogous to WER computation, producing matched word pairs for scoring. This requirement limits applicability to collections for which ground truth exists; a constraint that is not trivial for large-scale archival digitisation projects, where verified transcriptions may cover only a small sample of holdings.

A reference-free extension is methodologically feasible and represents a natural direction for future development. In this mode, the researcher would compute the marginal log-probability of each predicted token in context, equivalent to a per-token perplexity decomposition, without comparison to any reference. Tokens to which the scorer assigns anomalously low probability given their context would be surfaced as candidates for review. This approach would enable preliminary screening of unverified transcriptions, supporting triage and prioritisation of manual correction efforts in large-scale digitisation workflows. It cannot, however, distinguish between a genuinely erroneous token and a rare but correct one, nor can it quantify semantic distance from the correct form. The reference-free mode thus trades precision for coverage, and is most appropriately understood as a complementary screening tool rather than a replacement for reference-based evaluation.


Wouter: 
- smoothness bias; connect to congition, interface of the conversational agent; transkribus interface, not knowing wht model is being used -- 
- faithfullness -
- reviled/revived example
- literature review: ASR; other papers on transcirption -- 

Andy:
work on Failure modes 
run misnomer on large collections? 
*must have an evaluation step if using VLMs (arc of the paper), think of archive
review description of our method
Transkribus interface issue, 
line between traditional OCR tool and VLM, 
add to failure mode: interpretation/perspective

CFP deadline: August 14 AoE

\printbibliography

\end{document}




% WH: below goes all the boilerplate template info, for reference

% \section{Introduction} 

% Here is an example of the first section of the paper. All standard LaTeX
% formatting commands work as expected, such as \textit{italic},
% \textbf{bold}, \texttt{code}, and \textsc{small caps}. 

% You may modify this file by renaming, deleting, or adding sections of
% your own and substituting our instructional text with the text of your paper. Add
% references to \texttt{biblography.bib} as BibTeX entries. These can then be
% cited by using the format at the end of this
% sentence \cite{tettoni2024discoverability}. You can also cite multiple papers
% together using the format at the end of this sentence
% \cite{barré2024latent, levenson2024textual, bambaci2024steps}.

% \subsection{Details} \label{sec:intro_details}

% You may also include subsections if they help organize your text, but they are not required. Use as many sections and subsections with whatever names work for your submission.

% \section{Elements}

% \subsection{Tables}

% Tables can also be added to the document using the standard LaTeX table
% format. Each table needs a unique label and caption. Below (in the source code)
% is an example of the code to create a table, along with a brief caption.

% \begin{table}[h]
%   \centering 
%   \begin{tabular}{cc}
%     \toprule
%     Column Name 1 & Column Name 2\\
%     \midrule
%     d1 & d2 \\
%     d1 & d2 \\
%     d1 & d2 \\
%     \bottomrule
%   \end{tabular}
%   \caption{Example table and table caption.}
%   \label{tab:example}
% \end{table}

% The table can be referenced as Table~\ref{tab:example}.

% \subsection{Figures}

% Figures can also be added to the document. As with tables, each figure needs
% a unique label and caption. The format is shown in the lines below
% (in the source code). Figure files themselves should be included along with the
% submission.

% \begin{figure}[t!]
%   \centering
%   \includegraphics[width=0.4\linewidth]{640x480.png}
%   \caption{Example figure and figure caption.}
%   \label{fig:example}
% \end{figure}

% A figure can be cited as Figure~\ref{fig:example}.

% \subsection{Equations}

% We can include mathematical notations using LaTeX mathematical formatting,
% such as:
% \begin{align}
% f(y) &= x^2. \label{fig:squared}
% \end{align}

% The line number of the equation can be cited as Equation~\ref{fig:squared}.

% \subsection{Other References}

% Finally, you can also cite other sections or subsections of your paper using
% the tags that you have used at the end of each of the section titles:
% Section~\ref{sec:intro_details}.

% % Print the biblography at the end. Keep this line after the main
% % text of your paper, and before an appendix. 
% \printbibliography

% % You can include an appendix using the following command
% \appendix

% \section{First Appendix Section} \label{appdx:first}

% Optional appendix sections can be included after the references section. 

% \end{document}
